\documentclass[11pt,a4paper]{article}

\usepackage[utf8]{inputenc}
\usepackage[T1]{fontenc}
\usepackage{amsmath,amssymb}
\usepackage{booktabs}
\usepackage{graphicx}
\usepackage[margin=2.5cm]{geometry}
\usepackage{hyperref}
\usepackage{xcolor}
\usepackage{listings}

\lstset{
  basicstyle=\ttfamily\small,
  keywordstyle=\color{blue},
  commentstyle=\color{gray},
  showstringspaces=false,
  frame=single,
  language=Matlab % closest builtin highlighting for Julia snippets
}

\title{JosephProjectors.jl: a Julia port of the parallelproj\\
  Joseph 3D PET projectors with Apple Metal support}
\author{J.~J.~G\'omez-Cadenas}
\date{June 2026}

\begin{document}
\maketitle

\begin{abstract}
We describe a Julia implementation of the non--time-of-flight Joseph 3D
matched forward and back projectors of the \textsc{parallelproj} library
(Schramm \& Thielemans, 2023). Using KernelAbstractions.jl, a single kernel
source runs multithreaded on the CPU and on GPUs, including Apple Metal,
a backend not available in the original C/CUDA implementation. The port
reproduces the reference algorithm exactly: forward projections on Metal are
bitwise identical to the CPU path, the forward/back pair satisfies the
adjointness identity to $\sim 10^{-10}$ relative deviation, and an Apple M1
Max sustains 57 (36) million LORs per second in forward (back) projection of
a $128^3$ volume.
\end{abstract}

\section{Introduction}

Iterative PET reconstruction algorithms (MLEM, OSEM, primal--dual methods)
require repeated application of a forward projector $A$, which maps a voxel
image to line integrals along lines of response (LORs), and its adjoint
$A^{\mathsf T}$ (back projector). The \textsc{parallelproj}
framework~\cite{schramm2023} provides fast C/OpenMP and CUDA implementations
of Joseph's method~\cite{joseph1982} wrapped in a Python API. Its compute
core is compact: six kernel families sharing per-LOR ``worker'' functions
written once and compiled for both CPU and CUDA.

This note documents a Julia port of the non-TOF workers. Two goals motivated
it: (i)~a native Julia API for reconstruction research, and (ii)~GPU execution
on Apple silicon through the Metal framework, which the original library does
not target. The port relies on KernelAbstractions.jl, which compiles a single
backend-neutral kernel source to CPU threads, Metal, CUDA, and other GPU
backends.

\section{Joseph's method}

\subsection{Forward model}

For LOR $i$ with endpoints $\mathbf{x}_s^{(i)}$ and $\mathbf{x}_e^{(i)}$ the
forward projection is the line integral of the trilinearly interpolated image
$\tilde f$,
\begin{equation}
  p_i \;=\; \int_{\mathrm{LOR}_i} \tilde f(\mathbf{x})\,\mathrm{d}\ell .
\end{equation}
Joseph's method chooses the \emph{principal axis} $d$ of the ray --- the axis
with the largest squared direction cosine,
$\cos^2\theta_k = \Delta r_k^2 / \lVert \Delta\mathbf{r} \rVert^2$ with
$\Delta\mathbf{r} = \mathbf{x}_e - \mathbf{x}_s$ --- and evaluates the
integral as a sum over the voxel planes perpendicular to $d$,
\begin{equation}
  p_i \;\approx\; c \sum_{j=j_{\min}}^{j_{\max}}
      \mathcal{B}_d\!\left[f\right]\!\left(\mathbf{u}_j\right),
  \qquad
  c = \frac{\Delta_d}{\lvert\cos\theta_d\rvert},
  \label{eq:joseph}
\end{equation}
where $\mathcal{B}_d[f](\mathbf{u}_j)$ is the bilinear interpolation of the
image in plane $j$ at the point $\mathbf{u}_j$ where the ray crosses it,
$\Delta_d$ is the voxel size along $d$, and the correction factor $c$ converts
the plane-by-plane sum into a length-weighted integral. Because consecutive
crossing points differ by a constant increment, the inner loop reduces to two
fused multiply--adds plus one bilinear interpolation (four samples) per plane.

\subsection{Ray--volume intersection}

The traversal limits $j_{\min}, j_{\max}$ follow from a slab-method
intersection of the ray with the image bounding box. For each axis $k$ the
parametric entry/exit values are
\begin{equation}
  t_{1,2} = \bigl(b_{\min/\max,k} - x_{s,k}\bigr)/\Delta r_k ,
\end{equation}
clamped to $[0,1]$ and combined across axes by
$t_{\min} = \max_k \min(t_1,t_2)$, $t_{\max} = \min_k \max(t_1,t_2)$; the ray
misses the box iff $t_{\max} < t_{\min}$. Rays parallel to a slab give
$\Delta r_k = 0$ and IEEE-754 produces $t = \pm\infty$ or NaN ($0 \cdot
\infty$); the C reference handles this by using \texttt{fminf}/\texttt{fmaxf},
which \emph{ignore} NaN operands. This is a genuine portability trap for the
Julia port: \texttt{Base.min}/\texttt{max} \emph{propagate} NaN, so the port
defines explicit NaN-ignoring helpers,
\begin{lstlisting}
ieee_min(a, b) = ifelse(isnan(a), b, ifelse(isnan(b), a, min(a, b)))
\end{lstlisting}
which compile to branch-free \texttt{select} instructions on all backends.

\subsection{Matched adjoint}

The back projector applies the exact transpose of the forward weights: for
each plane it \emph{scatters} $c\,p_i$ into the same four voxels with the same
bilinear weights, instead of gathering from them. Distinct LORs write
concurrently to the same voxels, so the scatter uses atomic floating-point
additions. Matched $A/A^{\mathsf T}$ pairs guarantee the monotone convergence
properties of MLEM/OSEM and make the operators directly usable as
custom adjoints in automatic differentiation.

\section{Implementation}

\subsection{Single-source kernels}

The original library factors each projector into a per-LOR worker function
(C header), compiled twice: into an OpenMP loop and into a CUDA kernel. The
Julia port preserves this design with KernelAbstractions.jl:
\begin{lstlisting}
@kernel function fwd_kernel!(proj, @Const(xstart), @Const(xend),
                             @Const(img), org, voxsize, n)
    i = @index(Global)   # one work item per LOR
    ...
end
\end{lstlisting}
The wrapper selects the backend from the array type at call time
(\texttt{get\_backend(img)}): \texttt{Array} runs as a threaded CPU loop,
\texttt{MtlArray} is JIT-compiled to Metal IR, \texttt{CuArray} to PTX. The
package therefore depends only on KernelAbstractions; Metal.jl or CUDA.jl are
loaded by the \emph{user} environment, and their package extensions provide
the backend dispatch and the lowering of \texttt{@atomic} to native atomic
instructions.

\subsection{Metal-specific considerations}

\begin{itemize}
\item \textbf{Precision.} Apple GPUs do not implement IEEE binary64. The
  reference kernels are entirely \texttt{float}, so the port uses
  \texttt{Float32} throughout and loses nothing relative to the original.
\item \textbf{Atomics.} The back projector requires atomic \texttt{Float32}
  addition. Apple silicon exposes native floating-point atomics (Metal
  Shading Language~3); Julia's \texttt{@atomic} lowers to them through the
  \texttt{AtomixMetalExt} extension. A standalone risk-check kernel validated
  correctness before the port began.
\item \textbf{Special functions.} The (future) TOF kernels need
  $\operatorname{erf}$, which MSL does not provide. A kernel-local
  Abramowitz--Stegun 7.1.26 polynomial approximation was validated on Metal
  with maximum absolute error $4.6\times 10^{-7}$, ample for Gaussian TOF-bin
  weights.
\item \textbf{Indexing.} All in-kernel indices are \texttt{Int32}; dynamic
  tuple indexing is replaced by explicit three-way selects to keep the
  generated Metal IR branch-simple.
\end{itemize}

\subsection{Memory layout}

Julia arrays are column-major whereas the C reference is row-major. The port
keeps the image as an $(n_0, n_1, n_2)$ Julia array and addresses voxels
through one explicit linear-index helper, so the algorithm is
layout-agnostic; only the helper encodes the convention.

\section{Validation}

Four checks run on every available backend in the package test suite:
\begin{enumerate}
\item \textbf{Analytic line integrals.} Axis-aligned rays through a uniform
  unit image of $n$ voxels of size $\Delta$ must integrate to exactly
  $n\Delta$. Measured maximum relative error: $6\times 10^{-8}$
  (\texttt{Float32} machine epsilon).
\item \textbf{Adjointness.} For random $x$, $y$ and $2\times 10^5$ random
  LORs, $\langle Ax, y\rangle$ and $\langle x, A^{\mathsf T} y\rangle$
  (accumulated in \texttt{Float64} on the host) agree to a relative deviation
  of $1.1\times 10^{-10}$ (CPU) and $3.1\times 10^{-10}$ (Metal).
\item \textbf{Miss handling.} Rays that do not intersect the volume produce
  exact zeros in both directions.
\item \textbf{Cross-backend consistency.} Forward projections on Metal are
  \emph{bitwise identical} to the CPU result; back projections differ only
  through atomic accumulation order, bounded by $2.3\times 10^{-7}$ relative.
\end{enumerate}

\section{Performance}

Throughput on an Apple M1 Max (\texttt{Float32}, $128^3$ image, 2\,mm voxels,
$5\times 10^5$ random LORs through a 320\,mm sphere), Julia 1.12:

\begin{center}
\begin{tabular}{lrr}
\toprule
Backend & Forward (Mlors/s) & Back (Mlors/s) \\
\midrule
CPU, 8 threads (KernelAbstractions) & 19.9 & 3.3 \\
Metal GPU (M1 Max, 32 cores)        & 57.3 & 36.0 \\
\bottomrule
\end{tabular}
\end{center}

The Metal back projection benefits most ($\sim 11\times$ over the CPU), as the
atomic scatter parallelizes well on the GPU while serializing CPU threads.
These figures are of the same order as those reported for discrete CUDA GPUs
in~\cite{schramm2023}, and comfortably support interactive research-scale
reconstruction on a laptop.

\section{Outlook}

The TOF sinogram and listmode projectors are direct extensions: they reuse
the ray--cube intersection and bilinear scatter primitives and add Gaussian
TOF-bin weights, whose only nontrivial ingredient on Metal (the
$\operatorname{erf}$ approximation) is already validated. A scanner-geometry
layer and \texttt{ChainRulesCore} adjoints for deep-learning reconstruction
are planned on top of the projector core.

\begin{thebibliography}{9}
\bibitem{joseph1982} P.~M. Joseph,
  \emph{An improved algorithm for reprojecting rays through pixel images},
  IEEE Trans.\ Med.\ Imaging \textbf{1} (1982) 192--196.
\bibitem{schramm2023} G.~Schramm and K.~Thielemans,
  \emph{PARALLELPROJ --- an open-source framework for fast calculation of
  projections in tomography},
  Front.\ Nucl.\ Med.\ \textbf{3} (2024) 1324562,
  \href{https://doi.org/10.3389/fnume.2023.1324562}{doi:10.3389/fnume.2023.1324562}.
\end{thebibliography}

\end{document}
